\section{Quantitative acquired geometry}\label{sec:geometry}
The exact criterion separates implementability from estimates on the objective. We now re-prove a regular quantitative consequence, retaining the all-incoming and clock hypotheses explicitly. This section's curvature and integer-allocation mechanism is inherited from the preceding article\cite{prior}; the exact criterion above is not inferred from it.

All carriers in this section are one compact convex $K\subset\R^d$ with nonempty affine interior. Fix constants comparing the chosen norm $\|\cdot\|$ to the Euclidean norm $|\cdot|$. Suppose every $\Lambda_t(q,a)$ has density at most $\rho$ relative to the displayed $d$-dimensional affine volume, for \emph{all} $q,a,t$. Suppose $g$ is Lipschitz and $\kappa$-strongly concave in Euclidean coordinates, $\kappa>0$. Assume $J_t$ is uniformly TV-continuous and
\begin{equation}\label{eq:regular_stability}
 W_1(\Lambda_t(q,a),\Lambda_t(q',a))\le\beta_t\|q-q'\|,
 \quad 0<\beta_t\le B,\qquad
 \prod_{j=t}^{t+b-1}\beta_j\le\theta<1
\end{equation}
for every complete block of a fixed length $b$. The same experimental constants are used for every horizon. Zero coefficients can be enlarged to positive sufficient coefficients; no claim of sharpness for $\beta_t$ is made.

\begin{theorem}[Regular acquired-geometry transfer]\label{thm:regular}
Under these hypotheses there are constants $0<c\le C<\infty$, depending only on $K,d$, the norm comparison, $\rho,\kappa,\Lip(g),B,b,\theta$, such that, for every $n\ge1$ and $M\ge1$,
\begin{equation}\label{eq:regular_ext}
 cM^{-2/d}\le R_{\cp}(n,M)-B_n^*
 \le R_{\ext}(n,M)-B_n^*\le CM^{-2/d}.
\end{equation}
If report laws also have the common null sets of Proposition~\ref{prop:clock}, then the autonomous class is nonempty exactly when $M\ge n+1$ and
\begin{equation}\label{eq:regular_aut}
 c(M-n)^{-2/d}\le R_{\aut}(n,M)-B_n^*
 \le C(M-n)^{-2/d}.
\end{equation}
All legal adaptive acquisition policies are included. Without the block-product assumption the same construction gives, with $M\ge n+1$ in the autonomous expression,
\begin{align}\label{eq:weighted_upper}
 R_{\ext}(n,M)-B_n^*&\le C_0M^{-p}\sum_{t=1}^n a_t,\\
 R_{\aut}(n,M)-B_n^*&\le C_0(M-n)^{-p}
         \left(\sum_{t=1}^n a_t^{1/(1+p)}\right)^{1+p},\notag\\
 p&=2/d,\qquad a_t=\Lip(g)\prod_{j=t}^{n-1}\beta_j.\notag
\end{align}
The autonomous upper does not need clock-neutral reports; that hypothesis is used only for its matching lower.
\end{theorem}
For a finite measure $\lambda$ and a partition $\mathcal P$, set
\[
 \mathcal J_v(\lambda,\mathcal P)=\sum_{C\in\mathcal P}
 \lambda(C)v(\bar\lambda_C)-\int v\,d\lambda.
\]
Zero-mass cells are omitted. This is the Jensen objective of the corresponding barycentric contraction. The following estimate includes singular acquired measures, although a clean power law needs an additional mass estimate.

\begin{lemma}[Boundary-valid acquired-mass curvature estimate]\label{lem:curvature}
Let $K\subset\R^d$ be compact, convex, and have nonempty interior. If $v$ is concave and Euclidean $L$-Lipschitz on $K$, then $-v$ has a globally convex $L$-Lipschitz extension $f$. For every finite measure $\lambda$ on $K$ and grid side $h>0$,
\begin{equation}\label{eq:curvbound}
 \mathcal J_v(\lambda,\mathcal P_h)
 \le C_d h^{2-d}\sum_C\lambda(C)
 \mu_f(\overline B_{3\sqrt d h}(z_C)).
\end{equation}
If $\lambda$ is a probability law with density at most $\rho$, then for every integer $L_0\ge1$ there is a Borel partition with at most $L_0$ cells such that
\begin{equation}\label{eq:gridrate}
 \mathcal J_v(\lambda,\mathcal P)\le C_{K,d}\rho L L_0^{-2/d}.
\end{equation}
Neither an interior support margin nor differentiability of $v$ is required.
\end{lemma}
\begin{proof}
Put $f_0=-v$. At each $x\in\operatorname{int}K$ select all supporting slopes of $f_0$. Their Euclidean norms are at most $L$: directional difference quotients can be taken in every direction at an interior point. Take the supremum, on all $\R^d$, of the resulting supporting affine functions. This supremum is finite, convex, $L$-Lipschitz and agrees with $f_0$ on $\operatorname{int}K$; continuity gives equality on $K$. In particular, no assertion about arbitrary Lipschitz extensions preserving convexity is being used.

First suppose the extension is smooth. For a cell center $z$, subtract its tangent plane and call the resulting convex function $w$. Thus $w\ge0$ and $w(z)=0$. Put $R=2\sqrt d h$. For any $x$ in the cell, $B_{R/2}(x)\subset B_R(z)$. Writing $\operatorname{av}$ for the normalized spherical average, the submean inequality for $w$, radial monotonicity, and the divergence theorem give
\[
 w(x)\le \frac{C_d}{R^d}\int_{B_R(z)}w
 \le C_d\operatorname{av}_{\partial B_R(z)}w
 \le C_d R\operatorname{av}_{\partial B_R(z)}\partial_r w
 =C_dR^{2-d}\int_{B_R(z)}\Delta f.
\]
Here radial monotonicity gives $\int_{B_R}w\le (R/d)\int_{\partial B_R}w$.
The radial inequality is $w(z+R\omega)\le R\partial_r w(z+R\omega)$. This argument also applies in dimension one, with the boundary integral interpreted as the sum over two endpoints. If $b_C$ is the conditional barycenter, the affine tangent terms cancel and
\[
 \int_Cf\,d\lambda-\lambda(C)f(b_C)
 =\int_Cw\,d\lambda-\lambda(C)w(b_C)
 \le\lambda(C)\sup_Cw.
\]
Mollify $f$ on $\R^d$. The convex mollifications converge uniformly on bounded sets, and their positive Laplacian measures converge weakly to $\mu_f$. The limsup of the mass of the radius-$R$ ball is bounded by the mass of the closed ball of radius $3\sqrt d h$. Summing the finitely many relevant cells proves \eqref{eq:curvbound}.

A Lipschitz convex function has uniformly bounded Laplacian mass on any fixed bounded region: choose a smooth nonnegative cutoff $\chi$ that equals one there and use
\[
 \int\chi\,d\mu_f=-\int\nabla\chi\cdot\nabla f
 \le L\int|\nabla\chi|.
\]
For $h$ bounded above, the enlarged cell balls have bounded overlap, and $\lambda(C)\le\rho h^d$. Equation~\eqref{eq:curvbound} therefore gives $C_{K,d}\rho Lh^2$. Enclose $K$ in a fixed cube and subdivide each coordinate into $\lfloor L_0^{1/d}\rfloor$ intervals. There are at most $L_0$ cells. Since $\lfloor L_0^{1/d}\rfloor\ge L_0^{1/d}/2$, one has $h\le C_K L_0^{-1/d}$. This proves \eqref{eq:gridrate}, including $L_0=1$, by changing the fixed constant. The extension is used only in this proof, not as an observer state.
\end{proof}

\begin{lemma}[Integer allocation with a baseline label]\label{lem:integer}
Let $a_t\ge0$, $p>0$, and $m\ge1$. There are integers $L_t\ge1$ with $\sum_t L_t\le n+m-1$ and
\begin{equation}\label{eq:integer}
 \sum_t a_tL_t^{-p}\le 2^p m^{-p}
       \left(\sum_t a_t^{1/(1+p)}\right)^{1+p}.
\end{equation}
When all $a_t=0$ choose $L_t=1$. When $m=1$ the same choice is valid.
\end{lemma}
\begin{proof}
Otherwise put $w_t=a_t^{1/(1+p)}$, $A=\sum_tw_t$ and
$L_t=1+\lfloor(m-1)w_t/A\rfloor$. The total count is as stated. For $m=1$, $\sum_t a_t\le A^{1+p}$. For $m\ge2$ and $w_t>0$, $L_t\ge(m-1)w_t/A$ gives $\sum_ta_tL_t^{-p}\le A^{1+p}(m-1)^{-p}\le2^p A^{1+p}m^{-p}$. Zero-weight phases still have one real label.
\end{proof}

\begin{proof}[Proof of Theorem~\ref{thm:regular}]
At each phase the action set is common to all incoming states. The Wasserstein bound therefore implies
$\Lip(V_t)\le\beta_t\Lip(V_{t+1})$, by comparing the same action at two states before taking the finite minimum. Thus $\Lip(V_t)\le a_t$. Concavity was established independently from fixed-programme affine risks.

Starting from the preparation, choose a Bellman-minimizing action at every current label, partition the \emph{full incoming mixture} at the next cut, and use its cell barycenters. Lemma~\ref{lem:calibration} makes this a calibrated machine under its own law. Every pre-pooling law obeys the density bound, since the bound holds at all incoming mixture states and actions. Lemma~\ref{lem:curvature} with $L_t$ cells and the exact identity \eqref{eq:slack} give
$R-B_n^*\le C_0\sum_ta_t L_t^{-p}$. Choosing $L_t=M$ proves the external weighted upper. Choosing the integers of Lemma~\ref{lem:integer} with $m=M-n$ and using disjoint phase labels proves the autonomous weighted upper with total count $1+\sum_tL_t\le M$. No exact filter is executed outside these labels.

For $k=n-t$, block decomposition yields the explicit estimate
\[
 a_t\le\Lip(g)\max(1,B^{b-1})\theta^{\lfloor k/b\rfloor}.
\]
Hence both sums in \eqref{eq:weighted_upper} are uniformly bounded. This proves the two horizon-uniform uppers, including $m=1$.

For the lower, condition on the full history immediately before the final action/report. Every conditional final predictive law has density at most $\rho$, so its mixture under any randomized adaptive policy does too. If $S$ is this full-history state and a final label $Z$ takes at most $L$ values, its best forecast has conditional state $Q=\E[S\mid Z]$. Strong concavity implies
\[
 \E g(Q)-\E g(S)\ge\frac{\kappa}{2}\E|S-Q|^2.
\]
The $L$ balls of radius $r=(2\rho v_dL)^{-1/d}$ about the possible values of $Q$ cover mass at most $1/2$. Thus $\E|S-Q|^2\ge r^2/2$. This also covers a randomized label assignment: $Q$ still has at most $L$ values. Since the full-history risk of this policy is at least $B_n^*$, the same lower holds above $B_n^*$. Use $L=M$ for checkpoint and external observers. For an autonomous observer, Proposition~\ref{prop:clock} gives $L\le M-n$. This completes the proof. The lower concerns acquired mass under every admissible policy, not a packing of a reachable set.
\end{proof}

\begin{proposition}[Finite strata with unnormalized mass]\label{prop:strata}
Replace the density hypothesis by a finite report-readable decomposition of each $\Lambda_t(q,a)$ into submeasures supported on fixed compact convex affine sets $K_i\subset K$, of dimensions $d_i\le r$, $r\ge1$. Positive-dimensional submeasures have densities at most $\rho_i$ in their own affine volumes; zero-dimensional types are atoms. Suppose every full-history terminal law has an $r$-dimensional submeasure of actual mass at least $w>0$ and density at most $\rho_*$. Retain the task and block hypotheses on $K$. Then \eqref{eq:regular_ext}--\eqref{eq:regular_aut} hold with exponent $2/r$. The lower constant retains a factor
\begin{equation}\label{eq:stratum_mass}
 \kappa\,w^{1+2/r}\rho_*^{-2/r}.
\end{equation}
The autonomous lower again requires common report null sets.
\end{proposition}
\begin{proof}
For a phase allocation $L\ge2s$, where $s$ is the number of types, assign one cell to every atomic type and at least $\lfloor L/s\rfloor$ cells to each positive-dimensional type, reducing counts if necessary so the total is at most $L$. Restrict the curvature lemma to each $K_i$ and each unnormalized component. Since $d_i\le r$, its loss is at most a fixed multiple of $\Lip(V_t)L^{-2/r}$. Type is available from the current report, so the pair (type, cell) is an executable label rule. All incoming edges of that pair are pooled together. The resulting barycenter is in $K_i$. For $L<2s$, one global cell costs at most $\Lip(V_t)\operatorname{diam}(K)$; enlarge the constant by $(2s)^{2/r}$. Integer allocation and block summation proceed as before.

For the lower choose radius $r_0=(w/(2\rho_*v_r L))^{1/r}$. The union of $L$ ambient balls intersects the affine component in volume at most $Lv_r r_0^r$, and therefore covers at most $w/2$ of that submeasure. The remaining mass contributes at least $(w/2)r_0^2$ to squared distortion. Strong concavity gives \eqref{eq:stratum_mass}. No conditional normalization removes $w$, and no mass from other strata is subtracted.
\end{proof}

The ambient Euclidean chart in the curvature estimate is the finite-dimensional affine span of $K$; the chosen predictive norm is compared to its Euclidean norm with fixed constants. The curvature measure $\mu_f$ is distinct from the task's strong-concavity constant $\kappa$. In Theorem~\ref{thm:regular}, $B_n^*$ remains horizon- and task-dependent. Uniformity in $n$ concerns the excess constants, not a common stationary learning curve across different terminal tasks. Optimization over adaptive acquisitions is within the known-kernel, finite-action, phase-legal interface and its uniform hypotheses.
